\begin{frame}{역사: ``벨만''과 ``동적계획법''의 이름}
  \kb{리처드 벨만}(Richard Bellman, 1920$\sim$1984) --- 2차대전
  중 탄약 포대에 ``최적 포사거리 배분'' 문제를 풀던 미군 수학자.
  전후 물자배분·비행 경로 최적화로 \textbf{``동적계획법''이라는
  이름 자체를 만들었다}.
  \begin{itemize}
    \item 큰 문제 = \textbf{작은 부분 문제 + 재귀} 로 쪼개는
      재귀 구조로 풀어, 1957년 저서 \textit{Dynamic Programming
      and Its Applications}에서 정리.
    \item \kb{``동적''} = \textbf{시간}(시간에 따라 상태가
      바뀐다), \kb{``계획법''} = \textbf{계획}(미래 보상을
      계획적으로 최대화).
  \end{itemize}
  \vskip0.4em
  {\footnotesize ``재귀적 구조로 큰 문제를 쪼개는'' 이 패턴은
  \textbf{컴퓨터과학 전반}에서 반복: 강화학습(DQN, 정책기반)·
  컴퓨터그래픽스·자연어처리(최장 공통 부분열). 강화학습에서는
  ``미래의 \textbf{불확실한} 보상을 어떻게 지금 계산할 것인가''
  에 적용 --- 4.1$\sim$4.3절의 모든 알고리즘은 같은 패턴의 세
  가지 변주(기대/최적 $\times$ V/Q).}
\end{frame}
